\Needspace{20\baselineskip}
\section{例题与完整推导}
\label{l13:sec:examples}

\subsection{从输入划分到最终直方图}

设输入长度为 $N=16$，数据依次为
\begin{equation*}
x=(0,2,1,2,\;3,2,0,1,\;2,3,3,1,\;0,2,3,2).
\end{equation*}
使用两个线程块，每块四个线程，即 $G=2$、$T=4$、$P=8$，运行代码\ref{l13:lst:private-hist}。桶数仍为 $256$，但只有桶 $0$ 到 $3$ 可能非零。要求写出各线程负责的输入、每块的局部结果、全局结果，并计算原子操作次数。

线程块 $r$ 中线程 $t$ 的全局编号为 $g=4r+t$，处理索引 $g$ 与 $g+8$。表\ref{l13:tab:worked-mapping}逐一列出两个轮次，其中 $i\mapsto x_i$ 表示“输入索引映射到该位置的值”。

\begin{table}[htbp]
\centering
\caption{两块、每块四线程的完整输入分配}
\label{l13:tab:worked-mapping}
\begin{tabular}{cccc}
\toprule
块编号 $r$ & 块内线程 $t$ & 第一轮 $i\mapsto x_i$ & 第二轮 $i\mapsto x_i$ \\
\midrule
$0$ & $0$ & $0\mapsto0$ & $8\mapsto2$ \\
$0$ & $1$ & $1\mapsto2$ & $9\mapsto3$ \\
$0$ & $2$ & $2\mapsto1$ & $10\mapsto3$ \\
$0$ & $3$ & $3\mapsto2$ & $11\mapsto1$ \\
\midrule
$1$ & $0$ & $4\mapsto3$ & $12\mapsto0$ \\
$1$ & $1$ & $5\mapsto2$ & $13\mapsto2$ \\
$1$ & $2$ & $6\mapsto0$ & $14\mapsto3$ \\
$1$ & $3$ & $7\mapsto1$ & $15\mapsto2$ \\
\bottomrule
\end{tabular}
\end{table}

块 $0$ 处理的八个值为 $0,2,1,2,2,3,3,1$，所以桶 $0$ 到 $3$ 的局部计数为 $H^{(0)}=(1,2,3,2)$。块 $1$ 处理的值为 $3,2,0,1,0,2,3,2$，局部计数为 $H^{(1)}=(2,1,3,2)$。每个局部向量的元素之和都是 $8$，符合各块分配的输入数。两块分别完成第二次同步后，逐桶合并得到
\begin{equation}
H=H^{(0)}+H^{(1)}
=(1,2,3,2)+(2,1,3,2)
=(3,3,6,4).
\label{l13:eq:worked-hist}
\end{equation}
这里为了简洁只列出前四个桶，其余 $252$ 个桶为零。总和 $3+3+6+4=16$，且逐桶结果与直接计数相同。

直接算法需要 $16$ 次全局原子加。私有化算法需要 $16$ 次共享内存原子加，每块有 $4$ 个非零桶，所以全局合并只做 $4+4=8$ 次原子加。由于代码仍为每块清零 $256$ 个桶，局部清零总数却是 $2\times256=512$ 次；这个小例子说明，全局原子次数减少并不足以保证整体加速，私有化还需要足够工作量来摊薄固定开销。

如果输入只保留前 $14$ 个元素，索引 $14$ 和 $15$ 不再处理。块 $1$ 的线程 $2$、$3$ 在第二轮没有有效输入，但仍参与局部初始化和两次同步。最终前四桶为 $(3,3,5,3)$，总计数为 $14$。这一边界变化只减少两次输入贡献，不改变线程块各阶段的参与规则。

\subsection{由浮点吞吐率反推 L2 原子操作比例}

设一次原子操作在 L2 完成时耗时 $5\,\mathrm{ns}$，在 DRAM 完成时耗时 $120\,\mathrm{ns}$；一次浮点运算耗时 $1\,\mathrm{ns}$。每个线程执行 $5$ 次原子操作和 $100$ 次浮点运算，因此每次原子操作对应 $20$ 次浮点运算。给定有效浮点吞吐率为 $0.2424$ GFLOPS，其中 $1$ GFLOPS 表示每秒 $10^9$ 次浮点运算。求在 L2 完成的原子操作所占比例。

采用等效单条串行执行流的简化计时模型：一个线程所代表的工作量中，L2 原子操作、DRAM 原子操作和浮点运算的用时直接相加；给定吞吐率按这一等效工作流换算。令 $p\in[0,1]$ 表示原子操作在 L2 完成的平均比例，则一次原子操作的平均耗时为 $5p+120(1-p)$ 纳秒。一份线程工作量的平均耗时为
\begin{equation}
t(p)=5\bigl[5p+120(1-p)\bigr]+100
=700-575p\quad\mathrm{ns}.
\label{l13:eq:average-time}
\end{equation}
式中最外侧的 $5$ 是每线程原子操作次数，括号里的 $5$ 是一次 L2 原子操作的纳秒耗时；最后的 $100$ 是全部浮点运算的纳秒耗时。

将给定吞吐率换算为每秒完成的线程工作量：每份工作量含 $100$ 次浮点运算，所以速率为 $(0.2424\times10^9)/100=2.424\times10^6$ 份工作量每秒。取倒数得到平均用时
\begin{equation}
t_{\mathrm{obs}}
=\frac{100}{0.2424\times10^9}\ \mathrm{s}
\approx412.5413\,\mathrm{ns}.
\label{l13:eq:observed-time}
\end{equation}
在这里，一次浮点运算每纳秒等于 $1$ GFLOPS，也可以直接计算 $100/0.2424$，并将所得时间单位写为纳秒。

令式\eqref{l13:eq:average-time}等于式\eqref{l13:eq:observed-time}，可得
\begin{equation}
p=\frac{700-412.5413}{575}
\approx0.49993
\approx50\%.
\label{l13:eq:l2-ratio}
\end{equation}
给定吞吐率只有四位小数，因此答案写作约 $50\%$。若取 $p=1/2$ 反向核验，则平均原子操作耗时为 $(5+120)/2=62.5\,\mathrm{ns}$；每份工作量耗时 $5\times62.5+100=412.5\,\mathrm{ns}$，对应 $100/412.5=0.242424\ldots$ GFLOPS，四舍五入即为 $0.2424$。

还可以检查两个端点：若全部原子操作在 L2 完成，平均耗时为 $125\,\mathrm{ns}$，吞吐率为 $0.8$ GFLOPS；若全部在 DRAM 完成，平均耗时为 $700\,\mathrm{ns}$，吞吐率约为 $0.142857$ GFLOPS。给定值位于二者之间，与求得的 $p\in[0,1]$ 一致。$50\%$ 是许多线程或多次执行中原子操作的平均比例，不要求每个含五次原子操作的线程恰好有同样的 L2 次数。

这个反推依赖已明确的串行等效计时口径。只有“单线程内部三类操作不重叠”，还不足以把真实多线程设备的总吞吐直接换算成单线程延迟；线程之间可能重叠执行，也可能共享原子服务资源并排队。若用于真实内核测量，需要额外确定并行执行与吞吐统计的关系，才能由总吞吐唯一反推存储层级比例。

\subsection{实现检查与关键关系}

对直方图实现，可以围绕完整的贡献路径检查正确性：每个输入是否恰好被处理一次，桶下标是否有效，每次共享更新是否被保留，以及结果读取是否发生在相关更新完成之后。表\ref{l13:tab:diagnosis}列出四种具体错误及其后果。

\begin{table}[htbp]
\centering
\caption{典型实现错误及其修正依据}
\label{l13:tab:diagnosis}
\small
\begin{tabularx}{\textwidth}{YYY}
\toprule
错误写法或条件 & 产生的问题 & 修正依据 \\
\midrule
每块 $128$ 线程，仅用小于 $256$ 的线程编号清零同编号桶 & 只有桶 $0$ 到 $127$ 被清零，后半数组初值未知 & 按块内线程数跨步遍历全部桶 \\
共享数组上的累加改成普通加一 & 同块多个线程命中同一桶时丢失更新 & 局部冲突仍需共享原子加 \\
第二次同步前读取局部计数并合并 & 局部统计尚未结束，后续贡献可能遗漏 & 所有局部更新结束后再提交结果 \\
多个块用普通加法合并到同一公共桶 & 不同块的局部总数可能互相覆盖 & 使用全局原子加，或另设独占输出的归并阶段 \\
\bottomrule
\end{tabularx}
\end{table}

数学结果始终由 $H_j=\sum_i\mathbf1\{b(x_i)=j\}$ 定义。网格跨步划分通过 $i=g+kGT$ 保证输入覆盖；块级私有化通过 $H_j=\sum_rH_j^{(r)}$ 保持结果；跳过零桶时，全局更新次数为 $\sum_rU_r\le\min\{N,GB\}$。这些关系分别连接输入分配、统计正确性和性能开销。

合并访存改善同一轮输入读取，原子操作保护冲突的读—改—写，私有化把许多公共更新聚合为局部总数，线程粗化则用更长的局部工作摊薄初始化与合并成本。将这几个机制放回各自负责的环节，才能解释为什么一个实现既正确、又在特定数据规模和分布下更高效。
